% !TEX program = pdflatex
\documentclass[10pt,twocolumn]{article}

\usepackage[margin=0.85in]{geometry}
\usepackage{amsmath,amssymb,amsfonts}
\usepackage{booktabs}
\usepackage{enumitem}
\usepackage{hyperref}
\usepackage{xcolor}
\usepackage{graphicx}
\usepackage{multirow}
\usepackage{array}
\usepackage[T1]{fontenc}
\usepackage{lmodern}
\usepackage{microtype}
% \usepackage{algorithm}
% \usepackage{algpseudocode}
\usepackage{bm}
\usepackage{caption}
\usepackage{subcaption}
\usepackage[numbers,sort&compress]{natbib}

\hypersetup{colorlinks=true,linkcolor=blue!60!black,citecolor=green!50!black,urlcolor=blue!70!black}

% Compact lists
\setlist{nosep,leftmargin=*}

% Shorthand
\newcommand{\R}{\mathbb{R}}
\newcommand{\E}{\mathbb{E}}
\newcommand{\Var}{\mathrm{Var}}
\DeclareMathOperator*{\argmin}{arg\,min}
\DeclareMathOperator*{\argmax}{arg\,max}
\newcommand{\clip}{\mathrm{clip}}

\title{%
Monocular Depth Estimation as a Domain-Invariant\\
Representation for Vision-Based Quadruped Parkour
}
\author{
Theo Lemay \\
\textit{Undergraduate Thesis} \\
\texttt{theo@example.edu}
}
\date{April 2026}

\begin{document}
\maketitle

% ====================================================================
\begin{abstract}
% ====================================================================
Training legged robots for parkour in simulation requires vision, but
transferring vision policies to the real world is hindered by the
sim-to-real appearance gap.  We hypothesize that \emph{monocular depth
estimation} provides a domain-invariant intermediate representation:
because a pre-trained depth model maps diverse RGB appearances to a
canonical depth space, a student policy trained on estimated depth
should generalize across visual domains without photorealistic rendering
or explicit domain randomization.

We present a complete two-stage pipeline---privileged teacher via PPO,
then DAgger-based student distillation through a CNN--GRU depth
encoder---for Unitree Go2 parkour in the Genesis simulator.  Through
systematic experiments, we identify \emph{depth signal quality} as the
primary bottleneck, not network architecture: upgrading Depth~Anything~V2
from 25M to 100M parameters improves gap success by +49 percentage points,
and adding terrain textures further improves the hardest gaps by +72~pp.
Our best DA2-based student reaches 87\% gap success, within 7~pp of the
ground-truth depth upper bound.  We further compare against a
frozen-ResNet-18 baseline that operates directly on RGB.  Domain
invariance experiments under visual perturbations (texture, lighting,
color jitter) test the core hypothesis in simulation.
\end{abstract}


% ====================================================================
\section{Introduction}
\label{sec:intro}
% ====================================================================

Legged robots that traverse irregular terrain---stairs, hurdles,
gaps---must perceive geometry.  In simulation, privileged information
such as ground-truth heightmaps makes this straightforward, but
real robots lack such oracles.  Depth cameras are a common substitute,
yet many platforms (including the Unitree Go2) ship with RGB cameras
only.

The standard approach---training a student policy on rendered depth
from a simulated depth camera---sidesteps appearance but still requires
that the depth signal transfers from simulation to reality.  An
alternative line of work uses photorealistic rendering
\cite{yu2024lucidsim} or extensive domain randomization
\cite{tobin2017domain} to close the visual gap, but both are
computationally expensive and require domain-specific engineering.

We propose a simpler route: \emph{use a pre-trained monocular depth
estimator as a fixed intermediate representation.}  Foundation models
such as Depth~Anything~V2 (DA2) \cite{yang2024depth} have been trained
on millions of diverse real images with synthetic supervision, learning
a mapping from arbitrary RGB to a canonical depth space.  Because this
mapping is trained to be invariant to appearance, a student policy that
consumes DA2 depth should inherit that invariance for free.

\paragraph{Thesis statement.}
We hypothesize that pre-trained monocular depth estimation provides a
domain-invariant intermediate representation for vision-based quadruped
parkour.  Because the depth estimator maps diverse RGB appearances to a
canonical depth space, a student policy trained on estimated depth in
simulation should generalize across visual domains without photorealistic
rendering or explicit domain randomization.

\paragraph{Contributions.}
\begin{enumerate}
  \item A two-stage teacher--student pipeline from privileged scandot
        heightmaps to monocular depth for quadruped parkour.
  \item Systematic identification of depth signal quality---governed by
        model scale and terrain texture---as the key bottleneck.
  \item A domain invariance experiment comparing DA2-based and
        direct-RGB students under visual perturbations.
\end{enumerate}

\paragraph{Organization.}
Section~\ref{sec:related} reviews related work.
Section~\ref{sec:math} presents the mathematical foundations.
Section~\ref{sec:arch} describes the system architecture.
Section~\ref{sec:experiments} reports experimental results.
Section~\ref{sec:domain} presents the domain invariance experiment.
Section~\ref{sec:conclusion} concludes.


% ====================================================================
\section{Related Work}
\label{sec:related}
% ====================================================================

\paragraph{RL for legged locomotion.}
Modern quadruped controllers are typically learned via reinforcement
learning in simulation.  RMA \cite{kumar2021rma} introduced the
adaptation module paradigm, where a privileged teacher distills
environment-specific knowledge into a deployable student.
Lee~et~al.~\cite{lee2020learning} demonstrated sim-to-real transfer
for rough terrain without explicit state estimation.

\paragraph{Vision-based parkour.}
Extreme Parkour \cite{cheng2024extreme} trains a scandot-based teacher
and distills into a CNN--GRU student that processes depth images,
demonstrating agile behaviors (jumping, climbing) on a real Go1.
SoloParkour \cite{zhuang2024soloparkour} extends this with a
single end-to-end policy.  Robot Parkour Learning
\cite{zhuang2023robot} uses a similar teacher--student approach with
explicit depth processing.  All three systems assume access to a
depth camera at deployment.

\paragraph{Sim-to-real for vision.}
LucidSim \cite{yu2024lucidsim} uses generative models to produce
photorealistic training images, achieving sim-to-real transfer for
vision-based locomotion without real-world data.  Domain
randomization \cite{tobin2017domain} perturbs visual properties
(textures, lighting, colors) during training so the policy learns
invariance through exposure.  Both approaches require significant
computational overhead.

\paragraph{Monocular depth estimation.}
MiDaS \cite{ranftl2020towards} pioneered affine-invariant depth
estimation across diverse datasets.  Depth~Anything~V2 (DA2)
\cite{yang2024depth} builds on DINOv2 \cite{oquab2024dinov2}
vision transformers trained with self-supervised learning on 142M
images, achieving state-of-the-art zero-shot depth estimation.  DA2
outputs are affine-invariant: they predict relative depth up to an
unknown scale and shift, which we handle through online EMA
normalization (Section~\ref{sec:depth_norm}).


% ====================================================================
\section{Mathematical Foundations}
\label{sec:math}
% ====================================================================

\subsection{POMDP Formulation}

We model the parkour task as a Partially Observable Markov Decision
Process $(\mathcal{S}, \mathcal{A}, \mathcal{O}, T, R, \gamma)$.
At each timestep $t$, the full state $s_t \in \mathcal{S}$ includes
robot configuration, terrain geometry, and physics quantities.
The agent receives a partial observation $o_t \in \mathcal{O}$
and selects action $a_t \in \mathcal{A} = \R^{12}$ (one target angle
per joint).

\paragraph{Observation decomposition.}
Observations decompose into three groups:
\begin{itemize}
  \item \textbf{Proprioception} $o_\text{prop} \in \R^{53}$: base
    angular velocity~(3), projected gravity~(3), velocity
    commands~(3), heading commands~(2), additional commands~(2),
    joint positions~(12), joint velocities~(12), previous
    actions~(12), foot contacts~(4).
  \item \textbf{Exteroception} $o_\text{ext}$: either scandot
    heightmaps ($\R^{132}$ for the teacher), depth images
    ($\R^{58 \times 87}$), or RGB images ($\R^{3 \times 58 \times 87}$).
  \item \textbf{Privileged} $o_\text{priv}$: friction, restitution,
    and other environment parameters available only during teacher
    training.
\end{itemize}

\paragraph{Action space and PD control.}
Actions are target joint position offsets $a_t \in \R^{12}$ scaled
by 0.25~rad.  A PD controller converts these to torques at 250\,Hz
(decimation ratio 5, yielding 50\,Hz policy frequency):
\begin{equation}
  \tau = K_p \cdot (a_t^{\text{scaled}} - q) + K_d \cdot (-\dot{q}),
  \label{eq:pd}
\end{equation}
where $K_p = \text{diag}(35, 35, 35, \ldots)$ and $K_d = \text{diag}(0.5, 0.5, 0.5, \ldots)$
are per-joint gains, $q$ is the current joint position vector, and
$\dot{q}$ is the joint velocity vector.

\subsection{PPO with Generalized Advantage Estimation}
\label{sec:ppo}

The teacher policy $\pi_\theta$ is trained with Proximal Policy
Optimization (PPO) \cite{schulman2017proximal}.  PPO maximizes a
clipped surrogate objective:
\begin{equation}
  L^{\text{CLIP}} = \E_t\!\left[\min\!\left(
    r_t \hat{A}_t,\;
    \clip(r_t, 1{-}\epsilon, 1{+}\epsilon)\,\hat{A}_t
  \right)\right],
  \label{eq:ppo}
\end{equation}
where $r_t = \pi_\theta(a_t|o_t) / \pi_{\theta_{\text{old}}}(a_t|o_t)$
is the probability ratio and $\hat{A}_t$ is the advantage estimate.
The clipping parameter $\epsilon = 0.2$ prevents destructively large
policy updates.

\paragraph{Generalized Advantage Estimation (GAE).}
Advantages are computed recursively via GAE
\cite{schulman2016high}:
\begin{equation}
  \hat{A}_t = \sum_{l=0}^{T-t} (\gamma\lambda)^l \delta_{t+l},
  \quad
  \delta_t = r_t + \gamma V(s_{t+1}) - V(s_t),
  \label{eq:gae}
\end{equation}
where $\gamma = 0.998$ is the discount factor and $\lambda = 0.95$
controls the bias--variance tradeoff.  At $\lambda = 0$, GAE reduces
to the one-step TD error (low variance, high bias); at $\lambda = 1$,
it becomes Monte Carlo estimation (high variance, low bias).

\paragraph{Policy parameterization.}
The policy outputs a Gaussian distribution
$\pi_\theta(a|o) = \mathcal{N}(\mu_\theta(o), \sigma^2)$, where
$\sigma$ is a learned per-action-dimension standard deviation
(state-independent).

\paragraph{Hyperparameters.}
Table~\ref{tab:ppo_hypers} lists the PPO training configuration.

\begin{table}[h]
\centering
\caption{PPO hyperparameters for teacher training.}
\label{tab:ppo_hypers}
\small
\begin{tabular}{@{}ll@{}}
\toprule
Parameter & Value \\
\midrule
Discount $\gamma$ & 0.998 \\
GAE $\lambda$ & 0.95 \\
Clip $\epsilon$ & 0.2 \\
Learning rate & $1 \times 10^{-3}$ \\
Mini-batches & 4 \\
Epochs per iteration & 5 \\
Entropy coefficient & 0.0 \\
Steps per env per iter & 48 \\
\# environments & 4096 \\
Simulator & Genesis (GPU) \\
\bottomrule
\end{tabular}
\end{table}

\subsection{Reward Function}

The teacher is trained with a shaped reward function combining
progress incentives, sparse bonuses, and regularization penalties.
Table~\ref{tab:rewards} lists all terms with their scales.

\begin{table}[h]
\centering
\caption{Reward function terms and scales.}
\label{tab:rewards}
\small
\begin{tabular}{@{}lrl@{}}
\toprule
Term & Scale & Description \\
\midrule
\texttt{tracking\_goal\_vel} & 3.0 & Track goal velocity \\
\texttt{tracking\_goal\_heading} & 1.5 & Track heading \\
\texttt{progress\_along\_course} & 2.0 & Forward progress \\
\texttt{waypoint\_reached\_bonus} & 20.0 & Sparse bonus \\
\texttt{success\_bonus} & 40.0 & Course completion \\
\texttt{termination} & $-50.0$ & Early termination \\
\texttt{collision} & $-2.0$ & Body contact \\
\texttt{feet\_stumble} & $-1.0$ & Foot stumble \\
\texttt{feet\_edge} & $-2.0$ & Foot on edge \\
\texttt{torques} & $-2{\times}10^{-4}$ & Torque penalty \\
\texttt{dof\_power} & $-2{\times}10^{-4}$ & Power penalty \\
\texttt{dof\_acc} & $-2{\times}10^{-7}$ & Joint acceleration \\
\texttt{action\_rate} & $-0.01$ & Action smoothness \\
\texttt{action\_smoothness} & $-0.01$ & Second-order smooth \\
\texttt{dof\_pos\_limits} & $-2.0$ & Joint limit penalty \\
\texttt{torque\_limits} & $-0.2$ & Torque limit \\
\texttt{lin\_vel\_z} & $-0.5$ & Vertical velocity \\
\texttt{ang\_vel\_xy} & $-0.05$ & Roll/pitch rate \\
\texttt{orientation} & $-0.5$ & Non-upright penalty \\
\texttt{flat\_back} & $-3.0$ & Torso orientation \\
\bottomrule
\end{tabular}
\end{table}

\subsection{DAgger Distillation}
\label{sec:dagger}

The student policy is trained via Dataset Aggregation (DAgger)
\cite{ross2011reduction}.  Unlike behavioral cloning (BC), which
trains on a fixed dataset of expert demonstrations and suffers
$O(T^2 \varepsilon)$ compounding error over horizon $T$, DAgger
iteratively collects on-policy student rollouts labeled by the
teacher, achieving $O(T\varepsilon)$ regret.

In our setting, at each environment step the student observes
$o_\text{prop}$ and a depth image, produces an action
$a_\text{student}$, but the \emph{teacher's} action $a_\text{teacher}$
is executed in the environment.  The student is trained to minimize:
\begin{equation}
  \mathcal{L} =
    \underbrace{\|a_\text{teacher} - a_\text{student}\|_2}_{\text{action loss}}
    + \underbrace{\|y_\text{oracle} - \hat{y}\|_2}_{\text{yaw loss}},
  \label{eq:dagger_loss}
\end{equation}
where $\hat{y} \in \R^2$ is the student's predicted heading
command and $y_\text{oracle}$ is the ground-truth heading.
Both loss terms are weighted equally ($\alpha_\text{action} =
\alpha_\text{yaw} = 1.0$).

\paragraph{Windowed BPTT.}
The student uses a GRU for temporal aggregation.  Na\"ive
backpropagation through the full rollout ($N = 120$ steps) is
memory-prohibitive.  Instead, we partition each rollout into
$\lfloor N/W \rfloor$ windows of $W = 24$ steps and backpropagate
through each window independently, detaching the GRU hidden state
at window boundaries:
\begin{equation}
  \mathcal{L}_{\text{total}} = \frac{1}{\lfloor N/W \rfloor}
    \sum_{k=1}^{\lfloor N/W \rfloor} \mathcal{L}_k,
  \label{eq:bptt}
\end{equation}
where $\mathcal{L}_k$ is the loss over window $k$.  Gradients from
all windows are accumulated before a single optimizer step.  This
bounds peak memory to $O(W)$ while providing temporal gradient flow
within each window.

\paragraph{MTS yaw injection.}
During training, the student's predicted heading is compared against
the oracle.  When the prediction error
$\|\hat{y}_\text{scaled} - y_\text{oracle}\|_2 < \tau_\text{yaw}$
($\tau_\text{yaw} = 0.6$), the predicted heading replaces the
masked heading in the student's proprioceptive observation,
creating an implicit curriculum: the student must first learn accurate
heading prediction before its heading estimate is used for
action selection.

\subsection{Monocular Depth as Domain-Invariant Representation}
\label{sec:depth_theory}

\paragraph{Scale-shift ambiguity.}
Monocular depth estimation is inherently ill-posed: a single image
cannot determine absolute scale.  DA2 is trained with an
affine-invariant loss, producing predictions
$\hat{d} = \alpha \cdot d_{\text{true}} + \beta$ up to unknown
scale $\alpha$ and shift $\beta$ that vary per frame.

\paragraph{Domain invariance.}
We define a representation $f: \mathcal{I} \to \mathcal{Z}$ as
\emph{domain-invariant} if the variance of $f(I)$ across visual
domains $\mathcal{D}$ is much smaller than the variance of the
input:
\begin{equation}
  \Var_{\mathcal{D}}\!\left[f(I_\mathcal{D})\right]
  \ll
  \Var_{\mathcal{D}}\!\left[I_\mathcal{D}\right].
  \label{eq:domain_inv}
\end{equation}
The DA2 pipeline acts as an information bottleneck: it maps
high-dimensional, domain-specific RGB images to a low-dimensional
depth space that retains geometric structure while discarding
appearance.  A student policy that consumes only this depth output
should therefore be invariant to visual domain shifts by
construction.

In contrast, a student that directly processes RGB must learn its
own invariances from data, making it vulnerable to domain shifts
not seen during training.

\paragraph{EMA normalization.}
\label{sec:depth_norm}
To handle DA2's affine ambiguity, we normalize inferred depth using
online exponential moving average (EMA) percentile statistics.  Let
$q_{t}^{(\text{lo})}$ and $q_{t}^{(\text{hi})}$ denote the 2nd and
98th percentiles of the raw depth at step $t$:
\begin{align}
  \bar{q}^{(\text{lo})}_{t} &= (1 - \alpha)\,\bar{q}^{(\text{lo})}_{t-1} + \alpha\, q_{t}^{(\text{lo})}, \\
  \bar{q}^{(\text{hi})}_{t} &= (1 - \alpha)\,\bar{q}^{(\text{hi})}_{t-1} + \alpha\, q_{t}^{(\text{hi})}, \\
  d_{\text{norm}} &= \frac{d - \bar{q}^{(\text{lo})}}{\bar{q}^{(\text{hi})} - \bar{q}^{(\text{lo})}} - 0.5,
  \label{eq:ema_norm}
\end{align}
where $\alpha = 0.02$ and $d_{\text{norm}}$ is clamped to
$[-0.5, 0.5]$.  This maps the depth range to a fixed interval
regardless of DA2's per-frame scale and shift.


% ====================================================================
\section{System Architecture}
\label{sec:arch}
% ====================================================================

\subsection{Training Pipeline}

The system follows a two-stage pipeline (Figure~\ref{fig:pipeline}):
\begin{enumerate}
  \item \textbf{Stage 1 (Teacher).}  A privileged teacher policy is
    trained via PPO with access to ground-truth scandot heightmaps.
  \item \textbf{Stage 2 (Student).}  A student policy is trained via
    DAgger, replacing the privileged scandots with depth (GT, DA2, or
    RGB).
\end{enumerate}

\begin{figure}[t]
  \centering
  % \includegraphics[width=\columnwidth]{figures/pipeline.pdf}
  \fbox{\parbox{0.9\columnwidth}{\centering\vspace{2em}
    \textbf{[Pipeline diagram]}\\
    Teacher (scandots) $\to$ Student (depth/RGB)
  \vspace{2em}}}
  \caption{Two-stage training pipeline.  The teacher uses privileged
  scandot observations; the student replaces these with either GT depth,
  DA2-inferred depth, or direct RGB through a frozen ResNet-18.}
  \label{fig:pipeline}
\end{figure}

\subsection{Teacher Network}

The teacher (ActorCriticRMA) processes observations through three
parallel encoders:
\begin{itemize}
  \item \textbf{Scandot encoder}: $\R^{132} \to \R^{32}$ via MLP
    $[132 \to 128 \to 64 \to 32]$.
  \item \textbf{Privileged encoder}: $\R^{k_\text{priv}} \to \R^{32}$
    via MLP.
  \item \textbf{History encoder}: 1D convolution over past
    proprioceptive observations.
\end{itemize}
The concatenated features plus proprioception feed into the actor MLP
$[512 \to 256 \to 128 \to 12]$ with ELU activations.  A parallel
critic network estimates $V(s)$ for GAE.

\subsection{Student Network}

The student (ActorCriticParkourStudent) replaces the scandot encoder
with a recurrent depth encoder (Figure~\ref{fig:student_arch}):

\paragraph{Depth backbone.}
A CNN processes single-channel depth images $(1 \times 58 \times 87)$:
\begin{center}
\small
Conv2d(1,32,5) $\to$ MaxPool(2) $\to$ ELU \\
$\to$ Conv2d(32,64,3) $\to$ ELU \\
$\to$ Flatten $\to$ Linear($64{\times}25{\times}39$, 128) \\
$\to$ ELU $\to$ Linear(128, 32)
\end{center}

\paragraph{Combination MLP.}
The 32-dim backbone output is concatenated with the 53-dim
proprioceptive observation (with heading commands masked to zero) and
processed by a two-layer MLP: $[85 \to 128 \to 32]$.

\paragraph{GRU temporal encoder.}
A single-layer GRU with hidden dim 512 processes the 32-dim
combination features.  The GRU provides temporal context critical for
handling the DA2 update interval (every 5 steps) and accumulating
geometric information over time.

\paragraph{Output heads.}
The GRU output feeds a linear layer with Tanh activation producing
a 34-dim vector: 32-dim terrain latent $z$ and 2-dim heading
prediction $\hat{y}$ (scaled by 1.5).

\paragraph{Actor MLP.}
The actor MLP takes $[o_\text{prop}; z] \in \R^{85}$ and outputs
12 joint targets via the same architecture as the teacher actor:
$[512 \to 256 \to 128 \to 12]$ with ELU and Hardtanh($\pm$100)
output clipping.

\begin{figure}[t]
  \centering
  \fbox{\parbox{0.9\columnwidth}{\centering\vspace{2em}
    \textbf{[Student architecture diagram]}\\
    Depth CNN $\to$ Combo MLP $\to$ GRU $\to$ [latent + yaw]\\
    [proprio; latent] $\to$ Actor MLP $\to$ actions
  \vspace{2em}}}
  \caption{Student recurrent depth encoder architecture.  The CNN
  depth backbone and proprioception feed into a GRU that produces
  a terrain latent and heading estimate.}
  \label{fig:student_arch}
\end{figure}

\subsection{ResNet-18 RGB Baseline}
\label{sec:resnet}

To compare DA2-mediated depth with direct RGB processing, we train a
student variant that replaces the depth CNN backbone with a frozen
ResNet-18 \cite{he2016deep} truncated at \texttt{layer2}.  The
ResNet processes 3-channel RGB images $(3 \times 58 \times 87)$ and
produces a 32-dim feature vector via adaptive average pooling and a
trainable linear projection.  The frozen backbone contributes
$\sim$683K parameters (no gradients), while only the 4K-parameter
projection layer and the GRU/actor are trained.  This provides a
direct RGB baseline with minimal training overhead: 3.2\,s/iter and
16\,GB VRAM vs.\ 10.8\,s/iter and 28\,GB for the DA2 pipeline.

\subsection{Terrain and Curriculum}

The Genesis simulator generates parkour terrains with three obstacle
families (stairs, hurdles, gaps) at four difficulty levels per family.
Gap widths range from 0.24\,m (Row~0) to 0.50\,m (Row~3).  During
teacher training, a curriculum advances robots to harder rows upon
success.


% ====================================================================
\section{Experiments and Results}
\label{sec:experiments}
% ====================================================================

\subsection{Setup}

All experiments use the Genesis GPU simulator on 4$\times$ NVIDIA
L40S GPUs.  The Unitree Go2 quadruped has 12 actuated joints.
Evaluation uses forced terrain (all robots on the same obstacle type
and difficulty) with 128--256 episodes per condition.  Success is
defined as completing the full obstacle course.

\subsection{Teacher and GT-Depth Baselines}

Table~\ref{tab:baselines} shows the teacher and ground-truth depth
student performance on gap terrain across difficulty rows.

\begin{table}[h]
\centering
\caption{Teacher and GT-depth baselines (gap success rate, \%).}
\label{tab:baselines}
\small
\begin{tabular}{@{}lcccc@{}}
\toprule
& Row~0 & Row~1 & Row~2 & Row~3 \\
\midrule
Teacher (GT scandots) & 98.5 & 100.0 & 99.0 & 98.5 \\
GT-depth student & 94.5 & 94.5 & 89.1 & 87.5 \\
\bottomrule
\end{tabular}
\end{table}

The GT-depth student achieves $\geq$87\% on all rows, confirming
the CNN--GRU architecture can extract sufficient geometric
information from depth images.  The gap to the teacher (5--11~pp)
reflects the inherent information loss from scandots to a single
depth viewpoint.

A zero-depth ablation (replacing all depth pixels with zero)
drops success to 0--8\%, confirming the student relies on depth
for gap navigation rather than exploiting proprioceptive heuristics.

\subsection{DA2 Depth Estimation Students}
\label{sec:da2_results}

\paragraph{DA2-small ($\sim$25M params).}
Initial experiments with the smaller DA2 model plateaued at
$\sim$56\% gap success on Row~0, despite extensive
hyperparameter search across learning rates, curriculum schedules,
and normalization schemes (Table~\ref{tab:da2_small}).
Over 30 training runs established this as a consistent ceiling.

\begin{table}[h]
\centering
\caption{DA2-small experiments (best of 30+ runs).}
\label{tab:da2_small}
\small
\begin{tabular}{@{}lcc@{}}
\toprule
Configuration & Row~0 & Row~1 \\
\midrule
Best (diff-lr + cosine) & 56.3\% & 33.6\% \\
Fixed-range normalization & $\sim$45\% & --- \\
Calibrated linear & $\sim$45\% & --- \\
Frozen actor & $\sim$40\% & --- \\
\bottomrule
\end{tabular}
\end{table}

\paragraph{Root cause: depth signal quality.}
Analysis of DA2-small depth maps on untextured terrain revealed
the core issue.  Without visual texture, DA2 produces nearly flat
depth estimates (terrain--depth correlation $r = -0.10$), making
gaps almost invisible.  The gap-vs-ground pixel signal spans only
4.5\% of the normalized output range---too weak for reliable
detection through the CNN encoder.

\paragraph{DA2-base ($\sim$100M params).}
Upgrading to the four-times-larger DA2-base model immediately
improved performance.  Row~0 improved from 56\% to 68\% (+12~pp)
and Row~1 from 34\% to 83\% (+49~pp).  However, Row~3 remained
near zero (0.8\%), indicating that model scale alone does not solve
the hardest gaps.

\paragraph{Terrain texture.}
Adding a checkerboard texture to the terrain surface provides DA2
with visual features to estimate depth from, since monocular depth
estimation relies on texture gradients, perspective cues, and
learned priors.  Combined with DA2-base, this produced a
breakthrough:

\begin{table}[h]
\centering
\caption{Effect of DA2 model size and terrain texture on gap
success rate (\%).}
\label{tab:da2_main}
\small
\begin{tabular}{@{}lcccc@{}}
\toprule
Configuration & Row~0 & Row~1 & Row~2 & Row~3 \\
\midrule
DA2-small & 56.3 & 33.6 & --- & --- \\
DA2-base & 68.4 & 83.2 & 26.6 & 0.8 \\
\textbf{DA2-base + texture} & \textbf{87.5} & \textbf{93.0} & \textbf{87.1} & \textbf{80.5} \\
\midrule
GT-depth student & 94.5 & 94.5 & 89.1 & 87.5 \\
Gap to GT & $-$7.0 & $-$1.5 & $-$2.0 & $-$7.0 \\
\bottomrule
\end{tabular}
\end{table}

The DA2-base + texture student (checkpoint 7000) reaches 80--93\%
across all rows, within 7~pp of the GT-depth upper bound.  Row~3
improved from 0.8\% to 80.5\%---a +80~pp gain from texture alone.

\paragraph{Negative results.}
Several architectural modifications did not improve performance:
temporal EMA smoothing of depth frames hurt convergence;
Sobel gradient channels produced $\sim$40\% (below baseline);
alternative normalizations (calibrated linear, fixed range, mean-std)
performed comparably to EMA percentile.  This reinforces the finding
that depth \emph{signal quality} dominates architectural choices.

\subsection{ResNet-18 RGB Baseline}

The frozen-ResNet-18 student (Section~\ref{sec:resnet}) achieves
98\% overall parkour success (stairs + hurdles) but only 69\% on
gaps.  It converges in $\sim$4500 iterations with 3.2\,s/iter, making
it computationally efficient.  However, its gap performance
(69\% vs.\ 87\% for DA2-base+texture) suggests that direct RGB
features, while sufficient for stair and hurdle geometry, lack the
precision needed for narrow gap detection.

\subsection{Key Finding: Signal Quality Over Architecture}

Across all experiments, the dominant factor in student performance
was the quality of the depth signal reaching the encoder, not the
encoder architecture itself.  Evidence:
\begin{itemize}
  \item Architecture ablations (Sobel gradients, scandot prediction,
    GRU size, frozen actor) produced $<$15~pp variation.
  \item Depth quality changes (model size: +12--49~pp; texture:
    +72~pp on Row~3) produced order-of-magnitude larger effects.
\end{itemize}
This suggests that future work should prioritize depth estimation
quality (larger models, better training data) over student
architecture complexity.


% ====================================================================
\section{Domain Invariance Experiment}
\label{sec:domain}
% ====================================================================

The preceding results establish that DA2-based depth estimation
enables strong parkour performance.  We now test the core hypothesis:
does this pipeline provide domain invariance?

\subsection{Hypothesis}

\begin{itemize}
  \item $H_1$: The DA2-based student maintains performance under
    visual domain shifts (texture, lighting, color), because the depth
    estimator maps diverse appearances to a canonical depth space.
  \item $H_0$: DA2 and direct-RGB students degrade equally under
    visual perturbations.
\end{itemize}

\subsection{Visual Perturbations}

We apply perturbations at \emph{test time only} (the students were
trained under fixed visual conditions).  Perturbations span three
categories:

\paragraph{Texture.}
Replace the training checkerboard with: no texture (flat grey), solid
colors (red, blue), perlin noise, or brick pattern.

\paragraph{Lighting.}
Modify scene illumination: dim (30\% of default), bright (240\%),
or side-lit (lateral sun direction).

\paragraph{Color jitter.}
Apply image-space perturbations before DA2 or ResNet processing:
brightness scaling ($\pm30\%$, $\pm60\%$), per-channel color shifts
($\pm15\%$, $\pm30\%$).

\paragraph{Combined.}
Simultaneous texture + lighting + color jitter at two severity
levels (mild, extreme).

\subsection{Protocol}

Three students are evaluated under each perturbation:
\begin{itemize}
  \item \textbf{DA2-base+texture}: the primary DA2 student
    (Section~\ref{sec:da2_results}).
  \item \textbf{ResNet-RGB}: the frozen-ResNet-18 baseline
    (Section~\ref{sec:resnet}).
  \item \textbf{GT-depth}: invariant control (unaffected by visual
    changes).
\end{itemize}
Each condition is evaluated with 256 episodes on forced gap terrain
(Row~0).  We report success rate and the delta from baseline
(training visual conditions).

\subsection{Success Criteria}

A student is classified as:
\begin{itemize}
  \item \textbf{Domain-invariant}: $<$10~pp drop from baseline across
    all perturbations.
  \item \textbf{Domain-sensitive}: $>$30~pp drop at extreme
    perturbation.
\end{itemize}
The thesis is supported if the DA2 student is invariant \emph{and}
the RGB student is sensitive.

\subsection{Results}

\emph{[Results pending---to be filled after running the domain
invariance evaluation sweep.]}

% Placeholder table
\begin{table}[h]
\centering
\caption{Domain invariance results: gap Row~0 success rate (\%) under
visual perturbations.}
\label{tab:domain_inv}
\small
\begin{tabular}{@{}lccc@{}}
\toprule
Perturbation & DA2 & ResNet & GT-depth \\
\midrule
Baseline & --- & --- & --- \\
\midrule
No texture & --- & --- & --- \\
Solid red & --- & --- & --- \\
Noise texture & --- & --- & --- \\
\midrule
Dim lighting & --- & --- & --- \\
Bright lighting & --- & --- & --- \\
Side lighting & --- & --- & --- \\
\midrule
Brightness $\pm30\%$ & --- & --- & --- \\
Brightness $\pm60\%$ & --- & --- & --- \\
Color $\pm15\%$ & --- & --- & --- \\
Color $\pm30\%$ & --- & --- & --- \\
\midrule
Combined mild & --- & --- & --- \\
Combined extreme & --- & --- & --- \\
\bottomrule
\end{tabular}
\end{table}


% ====================================================================
\section{Discussion and Conclusion}
\label{sec:conclusion}
% ====================================================================

\subsection{Summary}

We presented a pipeline for vision-based quadruped parkour that uses
monocular depth estimation as an intermediate representation.  Our
systematic experiments revealed that depth signal quality---governed
by model scale and terrain texture---is the primary bottleneck for
student performance, not network architecture.  The best DA2-based
student achieves 87\% gap success, within 7~pp of the ground-truth
depth upper bound.

\subsection{Implications}

If the domain invariance results confirm $H_1$, the DA2 pipeline
eliminates the need for photorealistic rendering or domain
randomization for sim-to-real vision transfer.  The depth estimator,
trained once on diverse real-world data, provides a canonical bridge
between any visual domain and the geometric features needed for
locomotion.

\subsection{Limitations}

\begin{itemize}
  \item \textbf{No real-world deployment.}  All experiments are in
    simulation.  Real-world validation is needed to confirm
    domain invariance transfers beyond simulated visual shifts.
  \item \textbf{Computational cost.}  DA2-base (100M params) runs
    at $\sim$40\,ms/frame on an L40S GPU, which may be too slow for
    real-time deployment.  Distillation into a smaller network or
    quantization could address this.
  \item \textbf{Texture requirement.}  DA2 needs surface texture to
    produce useful depth estimates.  Truly featureless surfaces (e.g.,
    uniformly painted walls) may degrade performance.  However, most
    real-world surfaces have natural texture.
\end{itemize}

\subsection{Future Work}

\begin{itemize}
  \item \textbf{Real-world deployment} on the Unitree Go2 with
    onboard RGB camera.
  \item \textbf{Larger depth models}: DA2-large (335M params) or
    Video Depth Anything for temporal consistency.
  \item \textbf{End-to-end fine-tuning} of the depth estimator
    jointly with the student policy.
  \item \textbf{Multi-task evaluation} across diverse real-world
    environments (indoor, outdoor, varied lighting).
\end{itemize}


% ====================================================================
% References
% ====================================================================
\bibliographystyle{plainnat}
\begin{thebibliography}{20}

\bibitem[Cheng et~al.(2024)]{cheng2024extreme}
X.~Cheng, K.~Shi, A.~Agarwal, and D.~Pathak.
\newblock Extreme parkour with legged robots.
\newblock In \emph{IEEE ICRA}, 2024.

\bibitem[He et~al.(2016)]{he2016deep}
K.~He, X.~Zhang, S.~Ren, and J.~Sun.
\newblock Deep residual learning for image recognition.
\newblock In \emph{IEEE CVPR}, 2016.

\bibitem[Kumar et~al.(2021)]{kumar2021rma}
A.~Kumar, Z.~Fu, D.~Pathak, and J.~Malik.
\newblock {RMA}: Rapid motor adaptation for legged robots.
\newblock In \emph{RSS}, 2021.

\bibitem[Lee et~al.(2020)]{lee2020learning}
J.~Lee, J.~Hwangbo, L.~Wellhausen, V.~Koltun, and M.~Hutter.
\newblock Learning quadrupedal locomotion over challenging terrain.
\newblock \emph{Science Robotics}, 5(47), 2020.

\bibitem[Oquab et~al.(2024)]{oquab2024dinov2}
M.~Oquab, T.~Darcet, T.~Moutakanni, et~al.
\newblock {DINOv2}: Learning robust visual features without supervision.
\newblock \emph{TMLR}, 2024.

\bibitem[Ranftl et~al.(2020)]{ranftl2020towards}
R.~Ranftl, K.~Lasinger, D.~Hafner, K.~Schindler, and V.~Koltun.
\newblock Towards robust monocular depth estimation: Mixing datasets for
  zero-shot cross-dataset transfer.
\newblock \emph{IEEE TPAMI}, 2020.

\bibitem[Ross et~al.(2011)]{ross2011reduction}
S.~Ross, G.~Gordon, and D.~Bain.
\newblock A reduction of imitation learning and structured prediction to
  no-regret online learning.
\newblock In \emph{AISTATS}, 2011.

\bibitem[Schulman et~al.(2016)]{schulman2016high}
J.~Schulman, P.~Moritz, S.~Levine, M.~Jordan, and P.~Abbeel.
\newblock High-dimensional continuous control using generalized advantage
  estimation.
\newblock In \emph{ICLR}, 2016.

\bibitem[Schulman et~al.(2017)]{schulman2017proximal}
J.~Schulman, F.~Wolski, P.~Dhariwal, A.~Radford, and O.~Klimov.
\newblock Proximal policy optimization algorithms.
\newblock \emph{arXiv:1707.06347}, 2017.

\bibitem[Tobin et~al.(2017)]{tobin2017domain}
J.~Tobin, R.~Fong, A.~Ray, J.~Schneider, W.~Zaremba, and P.~Abbeel.
\newblock Domain randomization for transferring deep neural networks from
  simulation to the real world.
\newblock In \emph{IEEE/RSJ IROS}, 2017.

\bibitem[Yang et~al.(2024)]{yang2024depth}
L.~Yang, B.~Kang, Z.~Huang, et~al.
\newblock Depth {A}nything {V2}.
\newblock \emph{NeurIPS}, 2024.

\bibitem[Yu et~al.(2024)]{yu2024lucidsim}
W.~Yu, A.~Sharma, G.~Varma, and D.~Pathak.
\newblock {LucidSim}: Learning visual robot policies with generated experience.
\newblock \emph{arXiv:2405.20745}, 2024.

\bibitem[Zhuang et~al.(2023)]{zhuang2023robot}
Z.~Zhuang, Z.~Fu, J.~Wang, et~al.
\newblock Robot parkour learning.
\newblock In \emph{CoRL}, 2023.

\bibitem[Zhuang et~al.(2024)]{zhuang2024soloparkour}
Z.~Zhuang, Z.~Zhao, and H.~Zhao.
\newblock {SoloParkour}: Constrained reinforcement learning for visual
  locomotion from privileged experience of soloist.
\newblock \emph{arXiv}, 2024.

\bibitem[Cho et~al.(2014)]{cho2014learning}
K.~Cho, B.~van Merrienboer, C.~Gulcehre, D.~Bahdanau, F.~Bougares,
  H.~Schwenk, and Y.~Bengio.
\newblock Learning phrase representations using {RNN} encoder-decoder for
  statistical machine translation.
\newblock In \emph{EMNLP}, 2014.

\end{thebibliography}

\end{document}
